
\documentclass[a4paper,USenglish,cleveref,autoref,thm-restate]{lipics-v2021}

\hideLIPIcs

\bibliographystyle{plainurl}

\title{Report on Computing LCS w.r.t General \(\mathcal{E}\mathcal{L}\)-TBoxes}

\author{Daniel Hettinger}{Paderborn University, Germany}{dhetti2@campus.uni-paderborn.de}{}{}

\authorrunning{Daniel Hettinger}

\begin{CCSXML}
<ccs2012>
<concept>
<concept_id>10010147.10010178.10010187.10003797</concept_id>
<concept_desc>Computing methodologies~Description logics</concept_desc>
<concept_significance>500</concept_significance>
</concept>
</ccs2012>
\end{CCSXML}

\ccsdesc[500]{Computing methodologies~Description logics} %TODO Is this correct?

\keywords{Least common subsumer, Decsription Logics, \EL{}} %TODO add comma-separated list of keywords

\nolinenumbers

% My shortcuts :) 
\newcommand{\EL}{{\(\mathcal{EL}\)}}
\newcommand{\interpretation}{\mathcal{I}}
\newcommand{\I}{\mathcal{I}}
\newcommand{\J}{\mathcal{J}}
\newcommand{\TBox}{\mathcal{T}}
\newcommand{\interdomain}{\Delta^{\interpretation{}}}
\newcommand{\interfunc}[1]{{#1}^{\interpretation{}}}
\newcommand{\subsumes}[2]{#1 \sqsubseteq_{\TBox} #2 }
\newcommand{\lcs}[2]{\text{lcs}_{\TBox}(#1, #2)}
\newcommand{\subconcept}[1]{\textsc{sub}(#1)}
\newcommand{\canonical}[2]{\interpretation_{#1, #2}}
\newcommand{\simulates}[4]{(#1, #2) \lesssim (#3, #4) }
\newcommand{\bisimulates}[4]{(#1, #2) \backsimeq (#3, #4) }
\newcommand{\len}[1]{\lvert #1 \rvert}
\newcommand{\tail}[1]{\textsc{tail}(#1)}

\begin{document}

\maketitle

\begin{abstract}
    Lorem ipsum dolor sit amet, consetetur sadipscing elitr, sed diam nonumy eirmod tempor 
    invidunt ut labore et dolore magna aliquyam erat, sed diam voluptua. 
    % TODO write abstract. 
\end{abstract}

\section{Introduction}
\label{sec:introduction}

\begin{itemize}
    \item \textbf{What are Description Logics? What is \EL{} and what is it used for?} 
    Before going into the \textit{least common subsumer} and describing what it is useful 
    for, it makes sense to briefly cover what DL is and it's areas of application. The 
    original paper cites % TODO add citations. 
    \item \textbf{What is a reasoning service?} Introduce subsumption checking as a typical 
    reasoning service and give a brief example of what subsumption checking does. 
    \item \textbf{Motivation for the \textit{least common subsumer} (lcs)}. 
    \cite{DBLP:conf/ijcai/BaaderKM99} describes the ``bottom-up'' approach of building a 
    knowledge base (which uses lcs). Can use this as a motivation for this paper. 
    \item \textbf{What are the lcs and msc?} Provide a high-level definition of the lcs. 
    \item \textbf{Introduce an example.} I'll use the two TBoxes \(\mathcal{T}_1\) and 
    \(\mathcal{T}_2\) from the original paper. To show that the lcs doesn't have to exist, 
    can highlight that there is no lcs of Penicillin and Carbapenem for for \(\mathcal{T}_1\) 
    (and show why there can't be) and mention what the lcs would be for \(\mathcal{T}_2\).
\end{itemize}

\subsection{Contributions}
    The contribution of this report will be an example (which isn't to be found in the 
    original report) of the step-by-step computation the lcs of two concepts w.r.t 
    \(\mathcal{T}_2\) via the decision procedure outlined in Theorem 18.


\subsection{Related Work}
        \begin{itemize}
            \item Other papers that attempt to compute lcs are \cite{DBLP:conf/ecai/LutzPW10}, 
            which introduces an extenstion of \EL{} in which the lcs always exists and 
            \cite{DBLP:journals/japll/BaaderST07} and \cite{DBLP:conf/esws/PenalozaT11}, 
            which devise algorithms for \textit{approximating} the lcs up to some \textit{k}. 
            % TODO find other papers that references this paper for the Related Work section. 
        \end{itemize}

    
\section{Preliminaries}
\label{sec:prelims}

\subsection{Description Logics and \EL{}}

    \begin{itemize}
        \item Define \EL{}. 
        \item Define what an interpretation ($\interpretation$) is. What are the domain 
        ($\interdomain$) and interpretation function ($\interfunc{\cdot}$)? How does $\interfunc{\cdot}$ 
        extend to complex \EL{} concepts?
        \item What is a GCI? What does it mean for an interpretation $\interpretation$ to satisfy a GCI? 
        \item What is a TBox, and what does it mean for an interpretation $\interpretation$ to be a 
        model of a TBox $\TBox$?
        \item A small example of an interpretation that models $\TBox_1$. 
    \end{itemize}

\subsection{Subsumption and Least Common Subsumer}
    
    \begin{itemize}
        \item What is a reasoning task?
        \item What is subsumption? In other words, given some concepts 
        $C$ and $D$ and a TBox $\TBox$, what is $\subsumes{C}{D}$? Give a small example of
        subsumption. 
        \item How is $\lcs{C}{D}$ formally defined? Use definition 1 from the paper. 
        \item What is a common subsumer? Important to clarify that there is a set 
        of common subsumers, but only one unique lcs. Can use $\TBox_2$ to
        illustrate this.  
        \item What is the k-$\lcs{C}{D}$? For this, need to formally define role-depth (can use definition from \cite{Baader_Horrocks_Lutz_Sattler_2017}). 
    \end{itemize}

\subsection{Canonical Models}
    
    \begin{itemize}
        \item Give a formal definition of $\subconcept{C}$ given some concept $C$. Can cite \cite{Baader_Horrocks_Lutz_Sattler_2017}.
        Also provide a small example. 
        \item Can use definition 3 to formally define a canonical model. That is, 
        given some TBox $\TBox$ and concept $C$, how to get the canonical model 
        $\canonical{C}{\TBox}$. Can demonstrate this by building canonical models 
        $\canonical{P}{\TBox_1}$, $\canonical{C}{\TBox_1}$, $\canonical{P}{\TBox_2}$, and 
        $\canonical{C}{\TBox_2}$, which will also be used in later examples. 
    \end{itemize}

\subsection{Simulation Relations}
    
    \begin{itemize}
        \item Define what a simulation between two interpratations is 
        using definition 5 from the paper. 
        \item Define $\simulates{\I}{d}{\J}{e}$ and $\bisimulates{\I}{d}{\J}{e}$. 
        \item Give an example of simulation. Define two interpratations $\I$ and $\J$ 
        s.t. $\J$ simulates $\I$ but not the other way around. Explicitly define the simulation
        $\mathcal{S}$ for $\simulates{\I}{d}{\J}{e}$. Show where $\simulates{\J}{e}{\I}{d}$
        doesn't hold. 
        \item With canonical models and simulations covered, can introduce lemma 6 from the paper.
        These are some properties of canonical models and gives us another way to characterize 
        subsumption. If there is space, might make sense to add outlines for the proofs of these properties. 
        Should also elaborate on what the properties mean. 
    \end{itemize}

\subsection{Tree Unravelling}
    
    \begin{itemize}
        \item Define an unravelling $\I_d$ given some interpretation $\I$ and $d \in \interdomain$ using definition 8. 
        \item Given an element $\sigma$ in an unravelling $\I_d$, define the length $\len{\sigma}$ and $\tail{\sigma}$. 
        \item Define a finite subtree of an unravelling $\I_d^l$. 
        \item Using $\TBox_2$, give an example of $(\canonical{P}{\TBox_2})_{d_p}$. 
        \item Define the $l$-characteristic concept using defintion 9. Show how $(\canonical{P}{\TBox_2})_{d_p}$ can be used
        to compute the $2$-characteristic concept $X^l(\canonical{P}{\TBox_2}, d_p)$. 
    \end{itemize}
\section{Main Part}

    \textcolor{red}{TODO: Find a fitting title for this section and the subsections therein. This section should include a detailed description of the chosen paper's results and their analysis.}

    We currently use the LIPIcs class for our seminar reports. Here is a brief introduction on how to use it. For further details, check out \url{https://submission.dagstuhl.de/series/details/LIPIcs#author}.

    \begin{definition}
        \label{def:definition}
        Use definition enviornments to state definitions. Use a unique label for each definition.
    \end{definition}

    \begin{lemma}
        \label{lem:lemmas}
        Use lemma enviornments to state lemmas. Use a unique label for each lemma.
    \end{lemma}

    \begin{theorem}
    \label{thm:theorems}
        Use theorem enviornments to state theorems. Use a unique label for each theorem.
    \end{theorem}

    \begin{corollary}
        \label{cor:corollaries}
        Use corollary enviornments to state corollaries. Use a unique label for each corollary.
    \end{corollary}

    \begin{proof}
        Use proof enviornments to sate proofs.
    \end{proof}

    You can reference these blocks later: \cref{lem:lemmas}, \cref{thm:theorems}.

    There are different types of math enviornments. Use \(x+y\) or $x+y$ for inline equations,
    \[
        x+y
    \]
    for single line equations (without numbering) and 
    \begin{equation}
      x+y
    \end{equation}
    with numbering, and use  
    \begin{align}
        a &= x+y\\
        b &= z
    \end{align}
    for aligned multiline quations.
    If you do not need the numbering, use
    \begin{align*}
        a &= x+y\\
        b &= z
    \end{align*}
    instead.

    Use bibtex for your references. Add any references to the bibliography.bib. Use \cite{DBLP:journals/cacm/Knuth74} to cite a reference in your report. You can find bibtex entries for most papers on \url{dblp.uni-trier.de}. It suffices to use the condensed versions. If there are multiple versions, use the journal version preferably, then the conference version, then the preprint version.

\section{Conclusion and Future Work}

\textcolor{red}{TODO: Conclude the results you presented and present possibilities for future work.}

%%
%% Bibliography
%%

\bibliography{bibliography}

\end{document}
